\subsection{The integral type cut}\label{subsec:selected-cut}
\hypertarget{source-selected-cut}{}

A scalar cohomology length is counted with the multiplicity of the corresponding good character. To determine which characters contribute, we must realize the restriction to good types on the integral Hecke module at each of the two distinguished primes.

\begin{definition}[Type group and theta support at a localization]\label{def:localized-type-support}
Fix a point of a chosen chain, with its two distinguished depth-zero
primes \(g_0,g_1\). These are the primes defining the type conditions.
For each coefficient localization \(i\), let \(O_i\) be the coefficient
DVR, \(M_i\) the integral Hecke module and \(X_i\) its set of generic
Hecke characters. Write \(Z_i\subseteq X_i\) for the characters with
the prescribed theta summand and \(Z_i^+\subseteq Z_i\) for the good
ones, whose remaining type polynomials satisfy the tests of Definition~\ref{hyp:retained-type-tests}. The row group is
\[
 \Gamma=\prod_{k=0}^1\,\prod_{v\in\mathcal P_k}
              \mathrm{GL}_r(\mathbb F_{g_k}),\qquad h=|\Gamma|,
\]
where \(\mathcal P_k\) indexes the type places at \(g_k\) up to the
specified conjugacy. The same \(h\) is used at all localizations of
this point. It may vary along the family.
\end{definition}

\begin{lemma}[Integral central operators]\label{lem:integral-central}
Let \(e\) be a central idempotent of \(\mathbb C[\Gamma]\). For any
embedding of a characteristic-zero coefficient DVR \(O\) into \(\mathbb C\) whose image contains a primitive \(h\)-th root of unity, the induced map on group algebras satisfies
\[
 he\in\operatorname{im}\bigl(Z(O[\Gamma])\longrightarrow Z(\mathbb C[\Gamma])\bigr).
\]
The assertion also holds for a sum of products of central
idempotents.
\end{lemma}

\begin{proof}
For \(\gamma\in\Gamma\), let \(T\) be left multiplication by
\(\gamma^{-1}e\) on the regular representation. Its diagonal entries
on the group basis all equal the coefficient of \(\gamma\) in \(e\).
Since \(e\) is central, \((\gamma^{-1}e)^{h+1}=\gamma^{-1}e\).
The eigenvalues of \(T\) are therefore zero or powers of the
primitive root. Its trace is in the embedded ring \(O\). Taking
this trace for each \(\gamma\) gives the coefficients of \(he\).
Coefficientwise injectivity descends centrality. Products of
commuting central idempotents are central idempotents, and the
scaled lifts may be added in the center. In particular the
clearing factor for a product in the type cut is \(h\).
\end{proof}

We use the point data tied to the integral module. A complex
automorphic point \(P\) over a Hecke character \(a\) has a nonzero
form space \(U_P\) of the required tuple and lift index. Every
\(a\in X_i\) has at least one such point. At each \((k,v)\) the
irreducible finite group types have monic polynomial labels over
\(\mathbb F_{g_k}\). The central idempotent of a label is the sum
of the central idempotents of the types with that label. Its
action on the complex row space is the projection to those
isotypic components, by complete reducibility of finite group
representations in characteristic zero. Acting at every
\((k,v)\) gives the product of these commuting projections. The good
selector \(e_0\) is the sum of the products specified by the good
label tuples.

\begin{definition}[Support of the good type projector]\label{def:good-projector-support}
At a localization of Definition~\ref{def:localized-type-support} with \(Z_i^+\ne\varnothing\), the support of the projector \(e_0\) on the point spaces \(U_P\) is
\[
 S_i=\{a\in X_i: \text{some }P\text{ over }a
                    \text{ and }F\in U_P\text{ satisfy }e_0F\ne0\}.
\]
\end{definition}

\begin{proposition}[The good support of the type selector]
\label{prop:cut-selector}
For the integral center action and the depth-zero dictionary
on these point models, \(Z_i\cap S_i=Z_i^+\). There is an integral Hecke-equivariant endomorphism \(T:M_i\to M_i\). On the full eigenspace \(V_x\) at \(x\) it satisfies
\[
 T|_{V_x}=\begin{cases}h\,\mathrm{id}_{V_x},&x\in Z_i^+,\\0,&x\in X_i\setminus S_i.\end{cases}
\]
At each split type place \(w\) at \(g_k\), every point over
\(S_i\) has wild inertia with characteristic polynomial a power
of \(X-1\), zero monodromy, and each repeated eigenvalue of the tame generator has order
dividing \(g_k-1=\mathrm Nw-1\).
\end{proposition}

\begin{proof}
\proofstep{Check the projector on all points of the theta locus.}
When \(a\in Z_i\), the isobaric partition includes the specified
character block of rank \(t\). Removing it and undoing the common
twist gives the rank-\(r\) datum used for the type dictionary.
To see that the removal is available at any \(P\) over \(a\), start
with a point over \(a\) whose isobaric partition contains this block. The isobaric
partitions of points over \(a\) agree up to permutation. The
character block thus occurs at \(P\) with the same twisting
convention.

The parameter surjectivity in the dictionary gives
a type parameter with the same local data. The form space
\(U_P\) lies in the summand for that parameter. The series support
identity acts there by one for a good parameter and by zero
otherwise. The resulting test on the characteristic-zero fibre is
\[
 a\in Z_i^+\quad\Longleftrightarrow\quad\text{the parameter left after removal is good}.
\]
For \(a\in Z_i^+\), transport a partition with good remaining parameter to \(P\)
by a permutation. Conversely, if the parameter after removal is good,
the partition at \(P\) gives \(a\in Z_i^+\). This gives
\[
 e_0 F=F\quad(a\in Z_i^+),\qquad
 e_0 F=0\quad(a\in Z_i\setminus Z_i^+),\qquad F\in U_P .
\]
\proofstep{Descend from the point spaces to the integral Hecke module.}
The fibre calculation determines where the scalar should act. Its use in estimating lengths depends on realizing that scalar by an endomorphism before tensoring with the fraction field.

The integral action of the center of \(O_i[\Gamma]\)
commutes with the Hecke action on \(M_i\). Its comparison with
the action on the form spaces detects any scalar that acts on
all \(U_P\) over \(a\), as a scalar on the entire generic
\(a\)-eigenspace of \(M_i\). This is the tied eigenspace spanning
condition. The coefficient DVRs in this comparison already
contain the roots of unity in Lemma~\ref{lem:integral-central}.

The two action formulas give
\[
 Z_i\cap S_i=Z_i^+,
\]
with inclusion
of \(Z_i^+\) using nonzero form spaces. Lemma~\ref{lem:integral-central}
and the integral center action give an endomorphism of \(M_i\)
acting by \(h\) on the full eigenspaces over \(Z_i^+\). For
\(a\notin S_i\), the same endomorphism acts by zero on the
full eigenspace since all the form spaces there have zero
image. This is the operator used in the length comparison.

\proofstep{Read the inertia polynomial from one occurring type, then use it throughout the fibre.}
If \(e_0F\ne0\), some good tuple product has nonzero
value on \(F\). Each factor has nonzero value since the
factors commute. Expanding a factor as a sum of individual
type idempotents gives an occurring type with the indicated
label at that place. The depth-zero dictionary attaches to
that label \(P_v(X)\) the inertia polynomial in residue-field
labels
\[
                         P_v(X)(X-u_v)^t .
\]
Here \(P_v\) is squarefree with nonzero constant term and
avoids the two forbidden reciprocal labels from the theta
block. The dictionary gives zero monodromy and
\(u_v\in\mathbb F_{g_k}^{\times}\). These assertions include
the conjugate places, with the inverse labels on conjugation.
All roots are nonzero. A root of multiplicity at least two
must come from the repeated rank-\(t\) block and therefore
belongs to \(\mathbb F_{g_k}^{\times}\). The dictionary attaches
distinct continuous tame inertia characters to distinct
nonzero labels.

For two such characters \(\theta,\theta'\) and a tame generator
\(t\), agreement at \(t\) would give
\[
 (\theta/\theta')(t)=1\ \Longrightarrow\
 \overline{\langle t\rangle}\subseteq\ker(\theta/\theta').
\]
The kernel is closed by continuity and the cyclic subgroup is dense,
so the two characters would coincide. The dictionary is
multiplicative in the residue-field labels, so eigenvalues
of multiplicity at least two have order dividing \(g_k-1\).
At these completely split places \(g_k=\mathrm Nw\). Thus the local parameter satisfies the type condition. Writing \(m\) for the local rank and \(\lambda\) for any repeated tame eigenvalue, the three tests read
\[
 P_{\rm wild}(X)=(X-1)^m,\qquad N=0,\qquad \lambda^{\mathrm Nw-1}=1.
\]

This allows the repeated character block.

This argument used one point \(P\) over \(a\) with \(e_0F\ne0\) for some \(F\in U_P\).
At any other point above \(a\) the equality of the isobaric
partitions gives the same local parameter at each split
type place. Consequently this condition holds at every point of the cut, including points whose form spaces were unnecessary for deciding membership in \(S_i\). The cut comparison
therefore applies to the whole retained Hecke module.
\end{proof}

\subsubsection*{Compatibility with the integral module and rank}

Proposition~\ref{prop:cut-selector} identifies the generic eigensystems on \(S_i\). Their integral structure is the image of the original lattice under the same projection. Write \(V=V_{S_i}\oplus V_{S_i^c}\) for the rational
decomposition and \(M_{S_i}\) for the projected lattice.
Its exact presentation is
\[
  0\longrightarrow M_i\cap V_{S_i^c}
    \longrightarrow M_i\longrightarrow M_{S_i}\longrightarrow0.
\]
A central operator commuting with Hecke preserves each
eigenspace, so it preserves the kernel and descends integrally.
Writing \(p\) for the cut projection and \(V_a\) for the generic eigenspace at a retained character, one has
\[
 p|_{V_a}=\operatorname{id}_{V_a}.
\]
Hence the same point-space tests
detect its scalar action before and after taking the cut.

The row labels, the
inertia polynomials and the attached representations at the
retained points still come from the original place and tuple.
The finite paired-character equations restrict with them.
In particular the restriction uses the original norm choice,
with fewer local characters to test.

The reduced space still realizes the same tame and Hodge--Tate
labels at each retained character. At places over the coefficient prime use local--global compatibility with the specified inertia and monodromy condition at the type primes and the unramified parameter elsewhere. For the de Rham, potentially semistable representation at an unramified type, the criterion is
\[
 \mathrm{WD}(R)|_{I_v}=1,\quad N=0\quad\Longrightarrow\quad R|_{G_{E_v}}\text{ is crystalline}.
\]

At the
type primes the monodromy and inertia calculation is used in
the potentially crystalline comparison. The Hodge--Tate jumps
are those of the algebraic level-zero row. Their multiplicity
one implies that the semisimplified representation has distinct
absolutely irreducible constituents. Every summand of a de Rham representation is de Rham. At a fixed embedding \(\tau\), if an irreducible \(U\) occurred at least twice, any jump \(a\) of \(U\) would give
\[
 \dim\operatorname{gr}^a D_{\rm dR,\tau}(R^{\rm ss})\ge 2\dim\operatorname{gr}^a D_{\rm dR,\tau}(U)\ge2,
\]
contrary to simplicity of the labelled jumps. Here \(R\) denotes the attached representation. Thus the retained eigenspaces carry both the required inertia polynomials and the distinct labelled jumps used to identify their multiplicities.

\subsubsection*{Multiplicities at a twisted coefficient place}

Projection retains entire eigenspaces over the selected characters. We now compare their total multiplicity with the rank of the twisted theta image at a single place.

Fix a lift \(j\), a twisted coefficient place \(w\) and the
twist \(\mu=\mu(w)\) assigned to it. The row is the internal
direct sum \(M_j=\bigoplus_\sigma M_{j,\sigma}\).
Write \(\Theta_\mu\) for the theta map on this row and
\(\Theta_\mu^*\) for its adjoint. On \(M_{j,\sigma}\)
the operator \(\Theta_\mu^*\Theta_\mu\) is a positive
scalar. The twisted zero-space is the image of \(\Theta_\mu\)
and has dimension \(d_j(w)=\dim M_j\).

\begin{proposition}[Multiplicity after the type cut]
\label{prop:cut-multiplicity}
For a lift \(j\) and twisted coefficient place \(w\) on the cut
of Proposition~\ref{prop:cut-selector}, let
\(m_{i,\pi}\) be the dimension of the full good Hecke
eigenspace at a retained localization \(i\). Then
\begin{equation}\label{eq:retained-multiplicity}
 \sum_{i\ {\rm over}\ (j,w)}\;\sum_{\pi\in Z_i^+} m_{i,\pi}
             \le d_j(w)=d_j .
\end{equation}
\end{proposition}

\begin{proof}
\proofstep{Assign each good lower-row label by its integral split polynomials.}
For a localization \(i\) over \((j,w)\), assign a good
row label \(\sigma\) to \(i\) if the normalized polynomials
of its lifted parameter have integral representatives with
the residual Hecke values of \(i\) at every split place
away from the residue characteristic, the original
ramification set and the auxiliary level. Form the lifted
parameter by adding the rank-\(t\) character part and
taking the common normalizing twist. The local Hecke order
has a single residue character, so the test may use any
point of the localization. At surviving localizations it
may use any retained point. In the cut the polynomials,
their residue test and the assigned row label set are
therefore preserved. Projection of the split generic Hecke
algebra to a subset of characters retains their whole
eigenspaces. In particular, if \(V_a\) and \(V_a^{\rm cut}\) denote the original and cut eigenspaces at a good character, then
\[
 V_a^{\rm cut}\simeq V_a,\qquad \dim V_a^{\rm cut}=\dim V_a.
\]


\proofstep{Identify the full good eigenspace by the theta map and its adjoint.}
For a good character \(\pi\) and an isobaric row parameter
\(\Sigma\) for it, let \(M_{j,\Sigma}\) be the sum of row
spaces with that parameter and \(V_{i,\pi}\) the full
Hecke eigenspace. A row label belongs to this sum when its
local reading agrees with that of \(\Sigma\) at each split
place. We use the split local theta comparison
on the tied row and point models. The upward image of each
row type is in its specified point type. The local parameter
formula adds the character part and applies the common twist.
The split components of this base change determine the
Hecke character. The primes excluded from this split comparison have
Dirichlet density zero over \(E\).
For an inert place \(w\) above \(v\),
\[
 \mathrm Nw=(\mathrm Nv)^2.
\]

The Euler product of the lower field at \(2\) converges, so
\[
 \sum_{w\ {\rm inert}}(\mathrm Nw)^{-1}
   \le\sum_{v}(\mathrm Nv)^{-2}<\infty .
\]
The other exclusions form a finite set. In each finite Galois quotient
of \(G_E\), a nonempty conjugacy class has positive
Chebotarev density, so it meets the remaining Frobenius
classes. Local--global compatibility at those primes
and the common twist identify the functions
\[
       g\longmapsto \det(1-X\rho(g))
\]
for the two attached representations on a dense set.
Their continuity gives equality on \(G_E\). At all
unramified places the corresponding standard parameters
therefore agree. Strong multiplicity one for isobaric
representations identifies their constituents with multiplicity.
In the tied model these lists determine the local split
and unramified inert parameters used for the point
comparison.

Conversely, equality of the good split
Hecke readings gives equal split parameters by taking
the roots of the common Hecke polynomials and undoing
the normalization. \label{cor:selected-cut-1}
Thus \(\Theta_\mu M_{j,\Sigma}\subseteq V_{i,\pi}\).
For a row vector \(v\) in a different parameter and
\(u\in V_{i,\pi}\), use
\[
 \langle v,\Theta_\mu^*u\rangle
                    =\langle\Theta_\mu v,u\rangle .
\]
A nonzero value requires an occurring type over \(\pi\)
by orthogonality. Its split Hecke polynomials then agree
with those of \(\pi\), and cancelling the common character
block and twist identifies its row parameter with \(\Sigma\).
Thus \(\Theta_\mu^*V_{i,\pi}\subseteq M_{j,\Sigma}\).
The upward injection and the injection of the adjoint
on the full good eigenspace
(Proposition~\ref{prop:downward-on-good-space}) give
\[
 \dim M_{j,\Sigma}\le\dim V_{i,\pi}\le\dim M_{j,\Sigma},
 \qquad \Theta_\mu M_{j,\Sigma}=V_{i,\pi}.
\]

\proofstep{Identify the assignment with the sum of full good eigenspaces.}
A nonzero good row component assigned to \(i\) lifts
into a point type. Its lifted parameter has the split
polynomials of the residual test. The local point exhaustion
on this model puts the lift in the localization \(i\).
The two type-prime labels of its rank-\(r\) part are those
of the row, so the point is good. Conversely a good point
supplies its own integral Hecke polynomials for the residual
test. Any row of its parameter has good labels, since
base changes over this point have the same split components,
and the character part and twist cancel as above. These
two inclusions give the sum of good row components assigned
to \(i\) as the sum of the \(M_{j,\Sigma}\) for good
characters at \(i\). Put \(M^+\) for this assigned sum. Its image and dimension are
\[
 \Theta_\mu(M^+)=\bigoplus_{\pi\in Z_i^+}V_{i,\pi},\qquad \dim\Theta_\mu(M^+)=\sum_{\pi\in Z_i^+}m_{i,\pi}.
\]
For the upper bound an empty good point set contributes zero.

\proofstep{Separate the residual Hecke characters before summing dimensions.}
Summing over localizations gives a rank bound only if distinct localizations count disjoint lower-row labels.

Distinct localizations over \((j,w)\) have disjoint
assigned row sets. Their coefficient DVRs are identified
compatibly with the complex comparison. The same row
label has the same split parameter at both, and the added
character depends on \(j,\mu(w)\). If a label were
assigned to both localizations, choose a split place \(v\)
good for both and an index \(a\) whose Hecke generator has
different residues, using separation of the residual ideals.
The two lifted polynomials at \(v\) agree over the
identified DVR because their complex images agree.
Their coefficients of \(T^{r+t-a}\) have the same
reduction. Writing \(H_v(T)\) for the lifted polynomial and \(t_{v,a}\) for the Hecke value, the coefficient is
\[
 [T^{r+t-a}]H_v(T)=(-1)^a(\mathrm Nv)^{a(a-1)/2}t_{v,a}.
\]

The multiplier is a unit since \(v\) is away from the
residue characteristic of \(w\). Thus the two Hecke values
have the same reduction, a contradiction. Summing dimensions over the disjoint sets inside \(M_j\)
gives the inequality. The comparison of twisted and base
zero-spaces gives the last equality for this lift.
\label{cor:selected-cut-2}
\end{proof}

\subsubsection{The downward lift on the full good eigenspace}
\label{subsec:selected-point-adjoint}

The preceding multiplicity bound uses the rank of the whole good eigenspace. To compute it by the downward lift, we need injectivity of the adjoint on that eigenspace, including the multiplicity space of each occurring representation.

\begin{established}[Tied level, Weil data and good ranks]\label{hyp:tied-adjoint-level}
Consider the realizations at the fixed lift and its twist
on the selected cut of Proposition~\ref{prop:cut-selector}. Write \(A=V_j\) for the definite
Hermitian space used to realize the point coefficient type and
\(W=W_j\) for the lower space. Their dimensions and stable range are
\[
 (\dim_E A,\dim_E W)=(r+t,r),\qquad 0<r,\qquad t\ge r+5 .
\] The finite bad set \(S\) is the
set below the ramification and auxiliary level of this twist.
The source form has the fixed principal level and Cartan type,
and the splitting characters and Weil representation at both
levels are the ones in the tied realization. In particular
the host splits at the dyadic places in the application of
the first term identity.
\end{established}
\begin{proof}
Proposition~\ref{prop:cut-selector} selects points inside the fixed signed theta realization of Proposition~\ref{prop:signed-theta-data}. Projection to a type or an eigenspace restricts the vector space of functions and leaves its Hermitian spaces, splitting characters, Schwartz data and principal level fixed. The dimension at a lift is \(r+t\), with its lower constituent of dimension \(r\). The selected-rank estimates in Section~\ref{sec:family-window} give \(r>0\) and \(t\ge7r+4\ge r+5\). Finally \(-d\equiv1\pmod 8\) by the choice of \(d\). Hence \(\mathbb Q(\sqrt{-d})\otimes\mathbb Q_2=\mathbb Q_2\times\mathbb Q_2\). Base change to every dyadic completion of the totally real subfield splits the CM quadratic algebra there.
\end{proof}

\begin{proposition}[Injectivity on a full good Hecke eigenspace]
\label{prop:downward-on-good-space}
For the tied models of Lemma~\ref{hyp:tied-adjoint-level}
with \(r>0\) and \(t\ge r+5\), at each selected lift and twist let \(V\) be its full good Hecke eigenspace. The adjoint theta operator satisfies
\[
 \ker(\Theta_\mu^*|_V)=0.
\]
\end{proposition}

\begin{proof}
\proofstep{Read the good spherical system on both kinds of places.}
Let \(f\) be any element of the full Hecke eigenspace at a
good point. Read \(f\) in the vector-valued automorphic model
on \(A\) and take its inverse image \(u\) in the multiplicity
space. This reading is injective. The actions of the finite
Hecke operators on the reading and on the multiplicity space
are intertwined. In particular, if \(u\ne0\), then it has
the required spherical system of eigenvalues away from \(S\).
Here the assertion concerns a joint eigenvector, and the
system of eigenvalues can have arbitrary multiplicity.
The untwisted rank-\(r\) constituent and common normalizing
characters determined by the point give the following
property of the system.

For a good place \(w\) of \(E\)
over \(F\), let \(x_w\) be the unitary character value
of the twisted datum. The system contains the string
\[
 \left\{ x_w\,(\mathrm Nw)^{\,j-(t-1)/2}\mid 0\le j<t\right\}.
\]
Removing this multiset leaves \(r\) entries,
each of absolute value at most \((\mathrm Nw)^{(r-1)/2}\).
The local comparison with base change gives this description
at split and inert good places, including the passage from a
spherical parameter over \(F\) to the entries over \(E\).
The transported good projector acts as the identity on
\(u\). In proving this last assertion, use that the
archimedean signature and both selected finite type labels
of a localization form are good and that the arithmetic
idempotent reads as the same type projector on the
automorphic model. These two transports use the actual
twisted place and the same row index. \label{cor:selected-point-adjoint-1}

\proofstep{Detect a nonzero vector by the doubling integral.}
Suppose \(u\ne0\). At each real label take a vector \(v_\tau\)
of the prescribed lower compact type, and let \(w_0\) belong to
the definite source type. Use their Cartan inclusion
\[
 P=\iota_P((\bigotimes_\tau v_\tau)\otimes w_0)
\]
into the
Fock polynomials. The polarized vacuum test on the chosen
levels is nonzero on this Cartan summand. Pure tensors span it,
so choose \((v_\tau),w_0\) for which \(P\) times the vacuum
gives a nonzero vector \(\Phi\).

Take as \(a\) the scalar
component of the adelic extension of \(u\) paired against
this same \(w_0\). Schur orthogonality in the definite type gives
\[
 \int_{U(A)_\infty}
    \langle\omega(b)\Phi,\Phi\rangle\,r_a(b)\,db
  = c_1c_2\int_{U(A)_\infty}
       \big|\langle w_0,\rho(b)w_0\rangle\big|^2\,db ,
\]
with the conventions for the two coefficients chosen
conjugately. The Weil factor \(c_1\) is nonzero by the
nonzero chosen Cartan polynomial and the norm of its finite
part. The factor \(c_2\) is nonzero because the definite
multiplicity norm of \(u\) is positive. The remaining
continuous matrix coefficient is nonzero at identity.
This proves the real occurrence integral is nonzero and
also shows \(a\ne0\). Let \(p_0\) denote extraction of the chosen scalar component. At each good finite Hecke operator \(T\),
\[
 T a=T p_0(u)=p_0(Tu).
\]
Hence \(a\) has the same spherical parameter. The chosen finite Schwartz readings
are spherical away from \(S\) and principal-congruence
fixed at the level of \(a\). \label{cor:selected-point-adjoint-2}

The occurrence integral supplies the nonzero local factor in
the scalar zeta integral. To locate its pole, in the standard
partial \(L\)-function of \(a\) twist by
the inverse of \(x\). Write \(L_c^S\) for the Euler product on the \(r\) remaining entries after the twist. With \(w\mid S\) denoting all places over \(S\), the incomplete zeta factor is
\[
 \zeta_E^S(s)=\zeta_E(s)\prod_{w\mid S}(1-(\mathrm Nw)^{-s}).
\]
The product for the standard factor is
\[
 L_{\rm std}^S(s,a,x^{-1}) =
   L_c^S(s)\prod_{j=0}^{t-1}
          \zeta_E^S\big(s+(t-1)/2-j\big).
\]
The absolute parameter bound makes the Euler logarithm for
\(L_c^S\) normally convergent on
\(\RePart s>(r+1)/2\).

At \(s_*=(t+1)/2\) the last
zeta argument is \(1\), and all other arguments are integers
at least \(2\). The simple pole of the Dedekind zeta function
has nonzero residue, and none of its finitely many removed
factors vanish at \(1\). The factors with \(j<t-1\) are nonzero at \(s_*\), as is
\(L_c^S(s_*)\) by convergence of its Euler logarithm. Hence
\[
 \operatorname*{Res}_{s=s_*} L_{\rm std}^S(s,a,x^{-1})
  =L_c^S(s_*)\!
      \prod_{j<t-1}\zeta_E^S(t-j)
         \operatorname*{Res}_{u=1}\zeta_E^S(u)\ne0 .
\]
It is a simple pole.

Add \(t\) hyperbolic planes to \(W\) to obtain \(B\), of
rank \(r+2t\). Its doubled Siegel section is built from the
pair of copies of \(\Phi\boxtimes \Phi_H\), where \(\Phi_H\)
is the split Weil vacuum with principal-congruence support.
The complementary ranks in the doubled space and the positive doubling point satisfy
\[
 (r+2t)+r=2(r+t),\qquad s_0=\frac{(r+2t)-(r+t)}2=\frac t2.
\]

For clarity, the dimension parameter of the split group in
the first term identity is its Siegel parabolic rank \(r+t\).
Apply that identity in the complementary-space range,
with the compatible splitting character, the majorant
compact average, and the Ikeda map toward \(W\).
When \(W\) is anisotropic the paired residue is proportional
to the convergent integral of its theta kernel. When \(W\)
is isotropic it is the indicated leading regularized theta
coefficient on \(W\).

In this second case apply the
regularizing operator before integration, using the Witt
decomposition and the nonzero scalar of its auxiliary
Eisenstein series. Write \(\mathcal R\) for the paired residue before integration against the two forms. The implication used in both cases is
\[
 \Theta_{W,\phi}(a)=0\ \ (\text{all }\phi)\quad\Longrightarrow\quad \langle\mathcal R,a\otimes\overline a\rangle=0.
\]
Indeed, on doubled
Schwartz tensors the pulled-back theta kernel separates as
two theta factors, so integrating against the two forms
gives products of the corresponding lifts. The same
vanishing applies to the compact-averaged and regularized
inputs by the seesaw identity, continuity on Schwartz
vectors and the nonzero regularizing scalar.
\cite{gqt-rallis,yamana-first-term}\label{cor:selected-point-adjoint-3}

On the right of absolute convergence the doubling
unfolding instead gives the scalar zeta integral as
\[
 Z(s)=C\,\frac{L_{\rm std}^S(s+1/2,a,x^{-1})}{d_S(s)}.
\]
The denominator \(d_S\) is the unramified spherical
normalizer on the doubled group. For positive real part its
shifted unitary-character factors are absolutely
convergent and nonzero. The local factor \(C\) is independent
of \(s\) for the chosen data. To see this, the hyperbolic
vacuum coefficient vanishes outside the principal
congruence factor in the finite integration at level.

On this support the exponentiated height in the unfolding
is \(1\) and every right-level translate fixes the
coefficients. On the real compact group the hyperbolic
vacuum is fixed, leaving its nonzero squared norm times the
definite occurrence integral above. This yields \(C\ne0\).
Both \(C\) and \(d_S(s_0)\) are nonzero. The unfolding and
meromorphic continuation therefore give
\[
      \operatorname*{Res}_{s=s_0} Z(s)
        =\frac{C}{d_S(s_0)}
                 \operatorname*{Res}_{s=s_*}L_{\rm std}^S(s,a,x^{-1})
        \ne0,\qquad s_*=s_0+\tfrac12 .
\] This contradicts
vanishing of all the lifts and proves the existence of a
Schwartz vector and a scalar reading for which the integral
on \(W\) is nonzero. \label{cor:selected-point-adjoint-4}

\proofstep{Use the finite level generators to test the downward pairing.}
The test vectors used so far range over all Schwartz functions. At the selected level the adjoint uses a specified kernel; the finite level generators connect these two pairing tests.

The preceding conclusion uses all Schwartz data. To pass
to the actual downward operator, suppose its output on
\(u\) is zero. Its good target projection agrees with the
whole normalized level output on a good source by the
local type comparisons at the selected primes. The
output pairs the Cartan tensors against the coefficients
of the target compact types.

For a target argument \(g\) and a source linear functional \(\lambda\), denote the downward pairing with test \(\phi\) by \(B_{g,\lambda}(\phi)\). Vanishing on the generating tensors implies, by linearity,
\[
 B_{g,\lambda}(P\,\phi_{\rm vac})=0\qquad(\text{every allowed polynomial }P).
\]
Here \(\phi_{\rm vac}\) includes the fixed finite vacuum.
This uses the span of the real Fock polynomials under the
compact tensor maps and target translation. The same
vanishing holds for all Schwartz vectors whose finite
Heisenberg translates satisfy the linear relations of
the chosen finite vacuum. The real Schwartz approximation
on this fixed finite module is in the Sobolev norm for
which the downward pairing at each fixed argument is
continuous.

For a source level element \(k\) with character \(\chi_V(k)\),
right translation of the integration variable gives
\[
 I_u(\omega(k,1)\phi,l,g)=\chi_V(k)^{-1}I_u(\phi,l,g).
\]
Project a Schwartz test onto the inverse character by
averaging over the compact source level. On every fixed
finite vector this is a finite sum over cosets of an open
stabilizer after intersecting with the character kernel.
The difference from the projection pairs to zero. On the
inverse-character space let \(G\) be the finite level
quotient and \(k_b\) lifts of its elements \(b\). The resulting right action on the tests is
\[
 \pi(b)=\chi_{\rm flag}(k_b)^{-1}\omega(k_b^{-1},1),\qquad b\in G,
\]
where \(\chi_{\rm flag}\) is the flag part of the level character. The extra generators are the images of
\(\sum_{b\in G}c_b\pi(b)\) of the following types.

At one
of the selected residual places and for a nonzero vector
\(a\), take the subgroup of matrices fixing \(a^\perp\)
pointwise or inducing the identity on the quotient by
the line through \(a\). Use the identity at the other
places. For either choice of subgroup \(H\), write \(\eta\) for its prescribed determinant character or its inverse. The group algebra element acting on the source is
\[
 \sum_{b\in G}c_b\pi(b)=\sum_{b\in H}\eta(b)\pi(b).
\]

The selected good source projector kills these same
character sums with the normalized right action on
every scalar component. The substitution
\(h\mapsto hk_b\) in the theta integral transfers test
averaging to this source averaging, including the factor
\(\chi_{\rm flag}(k_b)^{-1}\). On a finite-dimensional orbit space \(V\), each subgroup \(H_j\) in this construction gives the normalized average
\[
 P_j=\frac1{|H_j|}\sum_{b\in H_j}c_b\pi(b),
\]
where the normalization is taken in characteristic zero. These projections
are self-adjoint in the invariant metric, so
\[
              \left(\sum_j\operatorname{im}P_j\right)^\perp
                            =\bigcap_j\ker P_j .
\]
In particular the self-adjoint projections give the orthogonal decomposition
\[
 V=\biggl(\bigcap_j\ker P_j\biggr)\oplus\biggl(\sum_j\operatorname{im}P_j\biggr).
\]
The good source
projector kills the pairings from the latter.

On the common kernel the local generation comparison
spans the inverse-character tests by target translates of
the fixed finite Heisenberg module above. This comparison
is applied one place at a time, with the spherical
generation at split and inert off-level places and the
prescribed local types at the twist and type primes.
Target translation replaces \(g\) by its right translate.
The normalized output thus tests a spanning family of the
inverse-character quotient and forces every downward
pairing of \(u\) to vanish. The pole test now applies in its full form:
\[
 \bigl(\Theta_{W,\phi}(u)=0\ \text{for every allowed }\phi\bigr)\quad\Longrightarrow\quad u=0.
\]

\label{cor:selected-point-adjoint-5}

\proofstep{Identify the normalized level operator with the actual adjoint.}
Finally identify this normalized level operator with the
point-to-row map \(\Theta_\mu^*\). The theta rule on simple
tensors is the adelic integral in the chosen source and
target coordinates. Both rules preserve level. They agree
after composition with the chosen embedding of the finite
coefficient model. The level comparison gives injectivity
of this restriction on the subspace of level rules, so the
resulting full coefficient tensors agree. Hence zero under
\(\Theta_\mu^*\) gives zero under the normalized level
operator. Applying its injectivity and the source coordinate isomorphism
proves \(f=0\) for every vector of the full eigenspace.
\label{cor:selected-point-adjoint-6}
\end{proof}

